% Taylor ve Maclaurin Serileri — teorik metin
% Derleme:  pdflatex -interaction=nonstopmode taylor-teori.tex   (iki kez)
\documentclass[11pt,a4paper]{article}

\usepackage[T1]{fontenc}
\usepackage[utf8]{inputenc}
% shorthands=off: Turkce babel ":", "!" ve "=" karakterlerini kisayol yapar;
% "=" aktif olunca enumitem'in leftmargin=* anahtari cozulemiyordu.
% Bize yalnizca Turkce heceleme lazim, kisayollar degil.
\usepackage[shorthands=off,turkish]{babel}
\usepackage{amsmath,amssymb,amsthm}
\usepackage{geometry}
\usepackage{graphicx}
\usepackage{xcolor}
\usepackage{enumitem}
\usepackage[hidelinks]{hyperref}

\geometry{margin=2.6cm}

\definecolor{vurgu}{HTML}{1B5E8C}
\definecolor{uyari}{HTML}{A03030}

\theoremstyle{plain}
\newtheorem{teorem}{Teorem}[section]
\newtheorem{onerme}[teorem]{Önerme}
\newtheorem{sonuc}[teorem]{Sonuç}
\newtheorem{lemma}[teorem]{Lemma}
\theoremstyle{definition}
\newtheorem{tanim}[teorem]{Tanım}
\newtheorem{ornek}[teorem]{Örnek}
\theoremstyle{remark}
% Ortam adi "not" TeX'in \not matematik komutuyla cakisiyordu;
% ortam acildiginda matematik kipi acilip tum belgeyi bozuyordu.
\newtheorem*{notkutusu}{Not}

\newcommand{\R}{\mathbb{R}}
\newcommand{\C}{\mathbb{C}}
\newcommand{\N}{\mathbb{N}}

\title{\textbf{Taylor ve Maclaurin Serileri}\\[0.4em]
\large Katsayıların zorunluluğundan kuantizasyonun doğuşuna}
\author{Taylor Laboratuvarı}
\date{\today}

\begin{document}
\maketitle

\begin{abstract}
Bu metin Taylor açılımını bir hesap tekniği olarak değil, bir \emph{zorunluluk
zinciri} olarak ele alır. Katsayılardaki $n!$ çarpanının neden kaçınılmaz
olduğundan başlanır; Taylor polinomu ile Taylor serisi arasındaki kritik ayrım
Lagrange kalanı üzerinden kurulur; yakınsaklık yarıçapının reel eksende değil
kompleks düzlemde belirlendiği gösterilir; pürüzsüz ama analitik olmayan
fonksiyonlarla $C^\omega \subsetneq C^\infty$ içermesinin gerçek olduğu
kanıtlanır; son olarak fizikteki iki büyük sonucun --- her kararlı dengenin
harmonik oluşu ve kuantizasyon --- aslında aynı serinin iki ayrı sonucu olduğu
savunulur. Metindeki her sayısal iddia bu deponun \texttt{02-python/}
klasöründeki betiklerle üretilmiş ve doğrulanmıştır.
\end{abstract}

\tableofcontents
\bigskip

% =====================================================================
\section{Katsayılar: $n!$ bir süsleme değildir}
% =====================================================================

Taylor açılımının en çok ezberlenen, en az anlaşılan yeri paydadaki $n!$'dir.
Oysa o çarpan bir tercih değil, tek bir isteğin kaçınılmaz sonucudur.

$a$ noktası civarında $f$'i bir polinomla temsil etmek isteyelim. Polinomu
$(x-a)$'nın kuvvetleri cinsinden yazmak doğaldır:
\begin{equation}\label{eq:genel-polinom}
  P(x) = c_0 + c_1(x-a) + c_2(x-a)^2 + \cdots + c_N(x-a)^N .
\end{equation}

Tek bir şey isteyelim: \emph{$P$'nin $a$ noktasındaki türevleri $f$'inkilerle
aynı olsun.} Bu isteğin katsayıları tamamen belirlediğini görelim.

\begin{onerme}\label{onerme:katsayilar}
$P$ \eqref{eq:genel-polinom} biçiminde olsun. Her $0 \le n \le N$ için
\[
  P^{(n)}(a) = n!\,c_n .
\]
\end{onerme}

\begin{proof}
$(x-a)^k$ terimini $n$ kez türevleyelim. $k < n$ ise sonuç sıfırdır. $k > n$
ise sonuç $(x-a)^{k-n}$ çarpanı taşır ve $x=a$'da sıfırlanır. Yalnızca $k=n$
terimi hayatta kalır ve
\[
  \frac{d^n}{dx^n}(x-a)^n = n(n-1)\cdots 2\cdot 1 = n!
\]
verir. Dolayısıyla $P^{(n)}(a) = n!\,c_n$.
\end{proof}

\begin{sonuc}[Katsayıların zorunluluğu]
$P^{(n)}(a) = f^{(n)}(a)$ istenirse, katsayılar için \emph{tek} seçenek
\begin{equation}\label{eq:katsayi}
  c_n = \frac{f^{(n)}(a)}{n!}
\end{equation}
olur.
\end{sonuc}

Buradaki $n!$, türevleme işleminin $(x-a)^n$ terimine yapıştırdığı çarpandır;
$n!$'e bölmek onu \emph{geri almanın} tek yoludur. Formülü ``şöyle
tanımlıyoruz'' diye sunmak, aslında bir teoremi tanım kılığına sokmaktır.

\begin{tanim}[Taylor polinomu ve Taylor serisi]
$f$, $a$'da $N$ kez türevlenebilir olsun.
\[
  P_N(x) = \sum_{n=0}^{N} \frac{f^{(n)}(a)}{n!}(x-a)^n
\]
ifadesine $f$'in $a$ merkezli $N$.\ dereceden \emph{Taylor polinomu} denir.
$f$ sonsuz kez türevlenebiliyorsa
\[
  T(x) = \sum_{n=0}^{\infty} \frac{f^{(n)}(a)}{n!}(x-a)^n
\]
biçimsel serisine $f$'in $a$ merkezli \emph{Taylor serisi} ($a=0$ için
\emph{Maclaurin serisi}) denir.
\end{tanim}

\begin{notkutusu}
Bu iki nesne \textbf{aynı şey değildir} ve aralarındaki fark bu metnin
omurgasıdır. $P_N$ her zaman vardır ve sonludur. $T$ ise yakınsayabilir de
yakınsamayabilir de; yakınsasa bile $f$'e eşit olmayabilir. Bölüm~\ref{sec:kalan}
birincisini, Bölüm~\ref{sec:analitik-degil} ikincisini ele alır.
\end{notkutusu}

% =====================================================================
\section{Taylor teoremi ve Lagrange kalanı}\label{sec:kalan}
% =====================================================================

Taylor teoremi bir yaklaşım tavsiyesi değil, bir \emph{eşitliktir}.

\begin{teorem}[Taylor, Lagrange kalanıyla]\label{teorem:taylor}
$f$, $[a,x]$ (veya $[x,a]$) kapalı aralığında $N$ kez sürekli türevlenebilir
ve açık aralıkta $(N+1)$ kez türevlenebilir olsun. O hâlde $a$ ile $x$
\emph{arasında} bir $\xi$ vardır ki
\begin{equation}\label{eq:lagrange}
  f(x) = P_N(x) + \frac{f^{(N+1)}(\xi)}{(N+1)!}(x-a)^{N+1}.
\end{equation}
\end{teorem}

\begin{proof}
$x \neq a$ sabitlensin ve $M$ sayısı
\[
  f(x) = P_N(x) + M\,(x-a)^{N+1}
\]
olacak biçimde tanımlansın ($x \ne a$ olduğundan $M$ tektir). Yardımcı
fonksiyonu
\[
  g(t) = f(t) - P_N(t) - M\,(t-a)^{N+1}
\]
ile tanımlayalım. Önerme~\ref{onerme:katsayilar} gereği $P_N^{(n)}(a) =
f^{(n)}(a)$ olduğundan
\[
  g(a) = g'(a) = \cdots = g^{(N)}(a) = 0,
\]
ayrıca $M$'nin seçimi gereği $g(x) = 0$'dır.

$g(a) = g(x) = 0$ olduğundan Rolle teoremi $a$ ile $x$ arasında bir $\xi_1$
verir: $g'(\xi_1) = 0$. Şimdi $g'(a) = 0 = g'(\xi_1)$ olduğundan Rolle yine
uygulanır ve $a$ ile $\xi_1$ arasında $g''(\xi_2) = 0$ olan bir $\xi_2$
bulunur. Bu $N+1$ kez tekrarlanırsa $a$ ile $x$ arasında
\[
  g^{(N+1)}(\xi) = 0
\]
sağlayan bir $\xi$ elde edilir. Öte yandan $P_N$ derecesi $N$ olan bir polinom
olduğundan $P_N^{(N+1)} \equiv 0$ ve $\frac{d^{N+1}}{dt^{N+1}}(t-a)^{N+1} =
(N+1)!$ olduğundan
\[
  0 = g^{(N+1)}(\xi) = f^{(N+1)}(\xi) - (N+1)!\,M
  \quad\Longrightarrow\quad
  M = \frac{f^{(N+1)}(\xi)}{(N+1)!},
\]
ki bu tam olarak \eqref{eq:lagrange}'dir.
\end{proof}

\begin{sonuc}[Hata sınırı]\label{sonuc:sinir}
$\displaystyle\sup_{t \in [a,x]} \bigl|f^{(N+1)}(t)\bigr| = M_{N+1} < \infty$ ise
\begin{equation}\label{eq:sinir}
  \bigl|f(x) - P_N(x)\bigr| \;\le\; \frac{M_{N+1}}{(N+1)!}\,|x-a|^{N+1}.
\end{equation}
\end{sonuc}

\subsection*{Sayısal doğrulama ve bir uyarı}

Bu deponun \texttt{02-python/src/kalan\_dogrulama.py} betiği, sekiz fonksiyon
$\times$ $N=1,\dots,12$ $\times$ $1501$ nokta olmak üzere toplam $144\,096$
noktada \eqref{eq:sinir} eşitsizliğini sınar. Sonuç: \textbf{teyit edilen ihlal
sayısı sıfırdır}.

Ancak burada ciddi bir sayısal tuzak vardır ve dürüstlük gereği açıkça
yazılmalıdır. Çift duyarlıklı (float64) aritmetikle bakıldığında
$2232$ nokta sınırı \emph{aşıyor görünür}. Örneğin $f=\exp$, $a=1$, $N=12$,
$x=0{,}996$ için
\[
  \text{gerçek hata} \approx 2{,}93\times 10^{-41},
  \qquad
  \text{float64 ile ölçülen} \approx 4{,}44\times 10^{-16}.
\]
Ölçülen değer sınırdan $10^{25}$ kat büyüktür. Bu bir teorem ihlali değil,
\emph{ölçüm aletinin çözünürlüğüdür}: $\varepsilon\,|f(x)| \approx 6\times
10^{-16}$ mertebesinin altındaki bir farkı float64 göremez. Aynı noktalar
\texttt{mpmath} ile $60$ basamakta yeniden ölçüldüğünde oranların tamamı $1$'in
altına döner. Betik bu yüzden float64'e yalnızca \emph{aday} işaretletir;
hükmü yüksek hassasiyete bırakır.

\subsection*{$\xi$ gerçekten vardır}

Teorem~\ref{teorem:taylor} bir \emph{varlık} iddiasıdır. Betik bunu somutlar:
\eqref{eq:lagrange}'den $\xi$'nin sağlaması gereken değer
\[
  f^{(N+1)}(\xi) = \frac{\bigl(f(x)-P_N(x)\bigr)(N+1)!}{(x-a)^{N+1}}
\]
olarak tek başına belirlenir ve bu denklemin kökü $(a,x)$ aralığında kök bulma
ile aranır. Her denemede bulunur. Özellikle dikkat çekici olan şudur:
$f = \ln(1+x)$, $a=0$, $x=1{,}8$ noktası yakınsaklık yarıçapının
\emph{dışındadır} --- seri orada ıraksar --- ama sonlu $N$ için eşitlik hâlâ
geçerlidir ve $\xi \approx 0{,}1436$ bulunur.

Bu, Taylor polinomu ile Taylor serisi ayrımının en keskin örneğidir:
\textbf{polinom her yerde anlamlıdır, seri her yerde anlamlı değildir.}

% =====================================================================
\section{Yakınsaklık yarıçapı kompleks düzlemin geometrisidir}\label{sec:yaricap}
% =====================================================================

\begin{teorem}[Cauchy--Hadamard]\label{teorem:ch}
$\sum_{n\ge 0} a_n (z-a)^n$ kuvvet serisi için
\[
  \frac{1}{R} = \limsup_{n\to\infty} |a_n|^{1/n}
\]
ile tanımlanan $R \in [0,\infty]$ sayısı şunu sağlar: seri $|z-a| < R$ için
mutlak yakınsar, $|z-a| > R$ için ıraksar.
\end{teorem}

\begin{teorem}[Yarıçap = en yakın tekilliğe uzaklık]\label{teorem:tekillik}
$f$, $a$ merkezli bir açık diskte holomorf olsun. $f$'in $a$ merkezli Taylor
serisinin yakınsaklık yarıçapı
\[
  R = \operatorname{dist}\bigl(a,\ \Sigma\bigr)
\]
olur; burada $\Sigma$, $f$'in analitik olarak devam ettirilemediği noktalar
kümesidir.
\end{teorem}

\subsection*{$1/(1+x^2)$: reel eksende hiçbir ipucu yoktur}

\[
  f(x) = \frac{1}{1+x^2}
\]
fonksiyonu reel eksende kusursuzdur: her yerde tanımlı, sınırlı ($0 < f \le 1$),
sonsuz kez türevlenebilir. Buna rağmen Maclaurin serisi
\[
  1 - x^2 + x^4 - x^6 + \cdots
\]
yalnızca $|x| < 1$ için yakınsar. \emph{Reel eksende bu sınırın sebebini
gösteren hiçbir şey yoktur.}

Sebep bir boyut yukarıdadır. $1+z^2 = (z-i)(z+i)$ olduğundan $f$'in $z = \pm i$
noktalarında kutupları vardır ve bu noktalar orijine tam olarak $1$
uzaklıktadır. Teorem~\ref{teorem:tekillik} gereği $R = 1$.

\subsection*{Merkezi kaydırmak}

Yarıçap fonksiyonun değil, \emph{(fonksiyon, merkez) çiftinin} özelliğidir.
Aynı $f$ için:
\begin{center}
\begin{tabular}{lll}
\hline
merkez $a$ & en yakın tekillik & $R = |a - (\pm i)|$ \\
\hline
$0$   & $\pm i$ & $1$ \\
$1$   & $\pm i$ & $\sqrt{2} \approx 1{,}4142$ \\
$2$   & $\pm i$ & $\sqrt{5} \approx 2{,}2361$ \\
\hline
\end{tabular}
\end{center}

\texttt{02-python/src/kompleks\_harita.py} betiği bunu \emph{iddia etmez,
gösterir}: kompleks düzlemde
\[
  L_N(z) \;=\; \frac{1}{N}\log_{10}\bigl|f(z) - P_N(z)\bigr|
\]
büyüklüğü çizilir. $|f - P_N| \sim C\,(|z-a|/R)^{N+1}$ olduğundan
\[
  L_N(z) \;\xrightarrow[N\to\infty]{}\; \log_{10}\!\left(\frac{|z-a|}{R}\right),
\]
yani $L_N$'nin \emph{işareti} doğrudan yakınsaklığı söyler ve $L_N = 0$ düzeyi
tam olarak $|z-a| = R$ çemberidir. Haritada disk kimse çizmeden belirir;
veriden çıkan sınır ile teorik çember çakışır.

\begin{notkutusu}[Sonlu $N$'in dürüst payı]
Haritadan okunan yarıçaplar teorik değerin sistematik olarak \%5--8 altında
çıkar. Bu bir hata değil, yukarıdaki yakınsamanın ancak $N\to\infty$ limitinde
tam olmasının sonucudur. Aynı sebeple $\ln(1+z)$ için sınır eğrisi sonlu
$N$'de gözle görülür biçimde çember \emph{değildir}: o fonksiyonun katsayıları
$1/n$ gibi cebirsel azalır (kutuplu fonksiyonlardaki geometrik azalmanın
aksine), bu yüzden sınır daha yavaş sıkışır.
\end{notkutusu}

\subsection*{Cauchy--Hadamard'ın sayısal sınırları}

\texttt{02-python/src/yakinsaklik\_orani.py} $R$'yi yalnızca katsayı dizisine
bakarak kestirir --- fonksiyona veya tekilliklerine hiç bakmadan. İki bağımsız
yolun aynı sayıda buluşması Teorem~\ref{teorem:ch}'nin içeriğidir. Betiğin
ortaya çıkardığı üç dürüst sınır:

\begin{enumerate}[leftmargin=*]
\item \textbf{$R=\infty$ sayısal olarak görülemez.} $\exp$ için
      $|a_n|^{1/n} = (1/n!)^{1/n} \approx e/n$ (Stirling); $n=60$'ta bu hâlâ
      $\approx 0{,}045$, yani $R \approx 22$. Sonsuzu sonlu veriden okumak
      mümkün değildir; yalnızca ``büyüyor'' denebilir.
\item \textbf{Cebirsel katsayılar yavaştır.} $\ln(1+x)$ ve $\sqrt{1+x}$ için
      $n=60$'ta bile \%7--13 sapma kalır. $\limsup$'ın doğası budur.
\item \textbf{Oran testi her zaman çalışmaz.} $f = 1/(1+x^2)$, $a=1$ için
      C--H \%1{,}66 hatayla doğru cevabı verirken oran testi \%22{,}66 sapar.
      Sebep: $i$ ve $-i$, $a=1$'e \emph{eşit uzaklıkta ama farklı yönlerdedir};
      katsayılar salınır ve $|a_n/a_{n+1}|$ oranının limiti yoktur. Teoremin
      neden limit değil $\limsup$ ile kurulduğu tam olarak budur.
\end{enumerate}

% =====================================================================
\section{Pürüzsüz ama analitik değil: $C^\omega \subsetneq C^\infty$}
\label{sec:analitik-degil}
% =====================================================================

\begin{tanim}
$f$ bir açık kümede sonsuz kez türevlenebiliyorsa $f \in C^\infty$ denir.
Ayrıca her noktanın bir komşuluğunda kendi Taylor serisine eşitse $f \in
C^\omega$ ($f$ \emph{analitiktir}) denir.
\end{tanim}

Açıkça $C^\omega \subseteq C^\infty$. İçermenin \emph{gerçek} olduğunu, yani
eşitlik olmadığını gösteren klasik tanık şudur.

\begin{teorem}\label{teorem:e-1x2}
\[
  f(x) = \begin{cases} e^{-1/x^2}, & x \neq 0,\\ 0, & x = 0\end{cases}
\]
fonksiyonu $\R$ üzerinde $C^\infty$'dur ve her $n \ge 0$ için $f^{(n)}(0) = 0$
sağlanır. Dolayısıyla $f$'in Maclaurin serisi özdeş olarak sıfırdır ve $f$
hiçbir $0$ komşuluğunda kendi Taylor serisine eşit değildir; yani
$f \notin C^\omega$.
\end{teorem}

\begin{proof}[İspatın çatısı]
$x \neq 0$ için tümevarımla
\[
  f^{(n)}(x) = p_n\!\left(\tfrac1x\right) e^{-1/x^2}
\]
biçiminde olduğu gösterilir; burada $p_n$ bir polinomdur. Gerçekten $n=0$ için
$p_0 = 1$; türev alındığında
\[
  f^{(n+1)}(x) = \left[-\tfrac{1}{x^2}p_n'\!\left(\tfrac1x\right)
  + \tfrac{2}{x^3}p_n\!\left(\tfrac1x\right)\right] e^{-1/x^2},
\]
ve köşeli parantez yine $1/x$ cinsinden bir polinomdur.

$x \to 0$ limitinde $u = 1/x \to \pm\infty$ konulursa
\[
  \lim_{x\to 0} f^{(n)}(x) = \lim_{|u|\to\infty} p_n(u)\,e^{-u^2} = 0,
\]
çünkü $e^{-u^2}$ her polinomdan hızlı sıfıra gider. Buradan tümevarımla
$f^{(n)}(0) = 0$ (türev tanımındaki fark bölümü de aynı sebeple sıfıra gider)
ve $f^{(n)}$'nin $0$'da sürekli olduğu çıkar. Dolayısıyla $f \in C^\infty$ ve
tüm Maclaurin katsayıları sıfırdır.

Seri $\equiv 0$ olduğundan her $x$ için yakınsar (yakınsaklık yarıçapı
sonsuzdur), ama $x \neq 0$ için $f(x) > 0$'dır. Seri yakınsar, fonksiyona
yakınsamaz.
\end{proof}

\texttt{02-python/src/analitik\_degil.py} betiği ilk dokuz türevi
\emph{sembolik limit} ile hesaplayarak sıfır olduklarını doğrular. Sayısal
türev burada kanıt olarak kabul edilemez: $f(0{,}1) = e^{-100} \approx
3{,}7\times 10^{-44}$ zaten float64'ün göremeyeceği kadar küçüktür, ``sıfır
çıktı'' demek hiçbir şey göstermez.

\subsection*{Sebep yine kompleks düzlemdedir}

Reel eksende $f$ uslu uslu sıfıra gider. Ama $z = iy$ alınırsa
\[
  -\frac{1}{z^2} = -\frac{1}{(iy)^2} = +\frac{1}{y^2}
  \quad\Longrightarrow\quad
  e^{-1/z^2} = e^{1/y^2} \xrightarrow[y\to 0]{} \infty .
\]
Aynı $z=0$ noktasına reel eksenden yaklaşınca $0$, hayali eksenden yaklaşınca
$\infty$ gelir. Limit yoktur: $z=0$ bir \emph{esas tekilliktir}. Somut olarak
$|z| = 0{,}05$ için reel yönde $1{,}9\times 10^{-174}$, hayali yönde
$5{,}2\times 10^{+173}$.

Bu, Bölüm~\ref{sec:yaricap}'daki dersin en uç hâlidir. Orada kompleks
tekillik yarıçapı $1$'e düşürüyordu; burada seriyi tamamen işe yaramaz kılıyor.

% =====================================================================
\section{Uygulama: her kararlı denge harmonik osilatördür}
% =====================================================================

$V$ bir potansiyel ve $x_0$ bir \emph{kararlı denge} noktası olsun; yani
$V'(x_0) = 0$ ve $V''(x_0) > 0$. $u = x - x_0$ yazıp açalım:
\begin{equation}\label{eq:potansiyel}
  V(x_0+u) = \underbrace{V(x_0)}_{\text{sabit}}
  + \underbrace{V'(x_0)\,u}_{=\,0}
  + \tfrac12 V''(x_0)\,u^2
  + \tfrac16 V'''(x_0)\,u^3 + \cdots
\end{equation}

Birinci terim sabittir; kuvvet $-V'$ olduğundan harekete katkısı yoktur.
İkinci terim, $x_0$'ın denge noktası \emph{olmasının tanımı} gereği sıfırdır.
Geriye ilk hayatta kalan terim olarak $\tfrac12 V''(x_0)u^2$ kalır --- yani
yay potansiyeli, $k = V''(x_0)$ yay sabitiyle.

\begin{sonuc}
Küçük salınımlar için, potansiyelin biçimi ne olursa olsun, hareket
\[
  m\ddot{u} = -V''(x_0)\,u
\]
harmonik osilatör denklemine indirgenir.
\end{sonuc}

Bu bir benzetme ya da kolaylık değildir. Doğa harmonik osilatörü
``seçmiyor''; Taylor serisi başka seçenek bırakmıyor. Harmonik osilatörün
fizikte her yerde çıkmasının sebebi, bir tesadüf değil,
\eqref{eq:potansiyel}'deki ikinci terimin denge noktasında zorunlu olarak
ölmesidir.

\subsection*{Atılan terim nereye gitti? Sarkaç}

Sarkaç denklemi $\ddot\theta = -(g/L)\sin\theta$ tamdır ve doğrusal değildir.
``Küçük açı yaklaşımı'' $\sin\theta \approx \theta$, Taylor açılımının
\[
  \sin\theta = \theta - \frac{\theta^3}{6} + \frac{\theta^5}{120} - \cdots
\]
ilk terimde kesilmesidir. Atılan $-\theta^3/6$ yok olmaz; \emph{periyotta}
geri döner. Tam periyot eliptik integralle verilir ve genliğe göre açılırsa
\begin{equation}\label{eq:periyot}
  T = T_0\left(1 + \frac{\theta_0^2}{16}
  + \frac{11\,\theta_0^4}{3072} + \cdots\right),
  \qquad T_0 = 2\pi\sqrt{L/g}.
\end{equation}

Harmonik osilatörde periyot genlikten bağımsızdır (izokronizm); gerçek
sarkaçta değildir. \texttt{02-python/src/sarkac.py} betiği tam denklemi
\texttt{solve\_ivp} ile çözer ve periyodu olay tabanlı sıfır geçişinden ölçer.
Bulgular ($L = 1$~m, $g = 9{,}80665$~m/s$^2$, $T_0 = 2{,}006409$~s):

\begin{center}
\begin{tabular}{rrrr}
\hline
$\theta_0$ & ölçülen $T$ (s) & $T_0(1+\theta_0^2/16)$ & harmoniğe göre sapma \\
\hline
$10^\circ$  & $2{,}010236$ & $2{,}010229$ & \%$0{,}19$ \\
$45^\circ$  & $2{,}086612$ & $2{,}083763$ & \%$4{,}00$ \\
$90^\circ$  & $2{,}368246$ & $2{,}315823$ & \%$18{,}03$ \\
$150^\circ$ & $3{,}535702$ & $2{,}865891$ & \%$76{,}22$ \\
\hline
\end{tabular}
\end{center}

Sayısal ölçüm ile eliptik kapalı biçim arasındaki en büyük fark
$1{,}3\times 10^{-11}$\%'dir; yani çözücü ile analitik sonuç birbirini
doğrular ve tablodaki her sapma Taylor kesmesinden gelir.

Mertebe kontrolü de yapılmıştır: düzeltmelerin kalan hatası log--log eğiminden
okunduğunda düzeltmesiz $\theta_0^{2{,}00}$, $\theta_0^2/16$ ile
$\theta_0^{4{,}00}$, bir sonraki terimle $\theta_0^{6{,}01}$ çıkar. Her
düzeltme mertebeyi tam olarak $2$ artırır, çünkü $\theta_0 \to -\theta_0$
simetrisi tek kuvvetleri yasaklar. \textbf{Taylor kesmesinin hata mertebesi
bir tahmin değil, ölçülebilir bir öngörüdür.}

% =====================================================================
\section{Bonus: Frobenius serileri ve kuantizasyonun doğuşu}
% =====================================================================

Şimdiye kadar bilinen bir fonksiyonun serisi çıkarıldı. Bu bölümde yön
tersine döner: \emph{seri önce gelir}, fonksiyon ondan doğar.

\subsection*{Legendre denklemi}

\begin{equation}\label{eq:legendre}
  (1-x^2)y'' - 2xy' + \lambda y = 0
\end{equation}
denkleminde $y = \sum_n a_n x^n$ konulursa katsayılar için
\begin{equation}\label{eq:legendre-rekurans}
  a_{n+2} = \frac{n(n+1) - \lambda}{(n+1)(n+2)}\,a_n
\end{equation}
rekürans bağıntısı çıkar. Denklem $\lambda$ üzerinde \emph{hiçbir kısıt
koymaz}: her $\lambda$ için bir seri çözüm vardır.

Kısıt fizikten gelir. \eqref{eq:legendre} küresel koordinatlarda
$x = \cos\theta$ ile kurulur, dolayısıyla $x = \pm 1$ (kutuplar) fiziksel
noktalardır ve çözümün orada sonlu kalması gerekir. Büyük $n$ için
\eqref{eq:legendre-rekurans}'den $a_{n+2}/a_n \to 1$ olur; seri $x = \pm 1$'de
ıraksar. Bundan kaçmanın \emph{tek} yolu payın sıfırlanmasıdır:
\[
  n(n+1) - \lambda = 0
  \quad\Longleftrightarrow\quad
  \lambda = \ell(\ell+1), \quad \ell \in \N .
\]
Seri o mertebede kesilir ve geriye bir \emph{polinom} kalır: Legendre
polinomu.

\subsection*{Hermite denklemi}

\begin{equation}\label{eq:hermite}
  y'' - 2xy' + 2\lambda y = 0,
  \qquad
  a_{n+2} = \frac{2(n-\lambda)}{(n+1)(n+2)}\,a_n .
\end{equation}
Kuantum harmonik osilatörde dalga fonksiyonu $\psi = H(x)e^{-x^2/2}$'dir. $H$
seri olarak kalırsa büyük $x$'te $e^{x^2}$ gibi büyür ve $\psi \sim e^{+x^2/2}$
olur --- normalize edilemez. Kesilme şartı $\lambda = n \in \N$'dir.

\texttt{02-python/src/frobenius.py} betiği $x=4$'te bunu ölçer: $\lambda = 3$
için $\psi(4) \approx -1{,}3\times 10^{-2}$ (sönüyor), $\lambda = 2{,}5$ için
$\psi(4) \approx 1{,}4\times 10^{1}$ (patlıyor).

\begin{notkutusu}[İki dürüst uyarı]
\emph{Birincisi}, kesilmeyen Legendre serisinin $x=1$'deki ıraksaması
logaritmiktir, yani çok yavaştır; $400$ terimde bile kısmi toplam sonlu bir
sayı gibi görünür. Sayısal bir deney ıraksamayı ``gösteremez'', yalnızca hiç
durmadığını gösterebilir. Iraksadığını söyleyen cebirdir.

\emph{İkincisi}, kesilme kararı katsayılara sayısal eşik uygulanarak
verilemez. $\lambda = 2{,}5$ için Hermite katsayıları da hızla küçülür
(tıpkı $e^{x^2}$'nin katsayıları gibi) ve eşik yöntemi seriyi yanlışlıkla
``kesilmiş'' sayar. Karar rekürans payının cebirsel sıfırlanmasından
okunmalıdır.
\end{notkutusu}

\subsection*{Sonuç}

Kuantum mekaniği dersinde $E_n = \hbar\omega(n+\tfrac12)$ diye ezberlenen
formüldeki $n$, rekürans bağıntısının hangi indekste sıfırlandığıdır.
Enerji kuantumlanmıştır çünkü seri ancak tam sayılarda kesilir. Tam sayılar
denklemden değil, \emph{çözümün sonlu kalması şartından} gelir.

\textbf{Kuantizasyon fizikten değil, yakınsaklıktan doğar.} Fiziğin koyduğu
tek şart ``çözüm sonlu olsun''dur; gerisini seri halleder.

% =====================================================================
\section{Toparlama}
% =====================================================================

\begin{enumerate}[leftmargin=*]
\item Katsayılardaki $n!$, türevlemenin ürettiği çarpanı geri almanın tek
      yoludur; formül bir tanım değil, bir zorunluluktur.
\item Taylor \emph{polinomu} ile Taylor \emph{serisi} farklı nesnelerdir.
      Lagrange kalanı sonlu $N$ için bir eşitliktir ve yakınsaklık yarıçapının
      dışında bile geçerlidir.
\item Yakınsaklık yarıçapı reel eksende değil kompleks düzlemde belirlenir:
      $R$, merkezin en yakın tekilliğe uzaklığıdır. Yarıçap fonksiyonun değil,
      (fonksiyon, merkez) çiftinin özelliğidir.
\item Serinin yakınsaması fonksiyona yakınsaması demek değildir:
      $C^\omega \subsetneq C^\infty$ ve tanığı $e^{-1/x^2}$'dir.
\item Her kararlı denge harmonik osilatördür, çünkü denge noktasında birinci
      mertebe terim ölür ve ilk hayatta kalan terim kareseldir.
\item Kuantizasyon, serinin kesilme şartından doğar.
\end{enumerate}

Ortak bir çekirdek vardır: \emph{Taylor serisi bir hesap aracı değil, yerel
davranışın tam bir muhasebesidir.} Atılan her terimin faturası bir yerde
kesilir --- sarkacın periyodunda, yakınsaklık diskinin kenarında, ya da bir
kuantum sayısında.

% =====================================================================
\section*{Kaynakça}
\addcontentsline{toc}{section}{Kaynakça}
% =====================================================================

\begin{thebibliography}{99}

\bibitem{rudin}
W. Rudin, \emph{Principles of Mathematical Analysis}, 3.\ baskı,
McGraw--Hill, 1976. (Taylor teoremi: Teorem 5.15; kuvvet serileri: Bölüm 8.)

\bibitem{ahlfors}
L. V. Ahlfors, \emph{Complex Analysis}, 3.\ baskı, McGraw--Hill, 1979.
(Yakınsaklık yarıçapı ve tekillikler: Bölüm 2 ve 4.)

\bibitem{stein}
E. M. Stein, R. Shakarchi, \emph{Complex Analysis}, Princeton University
Press, 2003.

\bibitem{whittaker}
E. T. Whittaker, G. N. Watson, \emph{A Course of Modern Analysis}, 4.\ baskı,
Cambridge University Press, 1927. (Legendre ve Hermite denklemleri,
Frobenius yöntemi.)

\bibitem{arfken}
G. B. Arfken, H. J. Weber, F. E. Harris, \emph{Mathematical Methods for
Physicists}, 7.\ baskı, Academic Press, 2013. (Seri çözümler ve özel
fonksiyonlar.)

\bibitem{landau}
L. D. Landau, E. M. Lifshitz, \emph{Mechanics}, 3.\ baskı,
Butterworth--Heinemann, 1976. ($\S$11: küçük salınımlar; $\S$28:
anharmonik titreşimler.)

\bibitem{griffiths}
D. J. Griffiths, \emph{Introduction to Quantum Mechanics}, 2.\ baskı,
Pearson, 2005. (Bölüm 2.3: harmonik osilatörün analitik çözümü ve kesilme
şartı.)

\bibitem{depo}
Bu deponun sayısal laboratuvarı: \texttt{02-python/src/} altındaki betikler ve
\texttt{out/tablolar/} altındaki üretilmiş tablolar. Metindeki her sayı oradan
gelir.

\end{thebibliography}

\end{document}
